\documentclass[11pt]{article}

% =====================================================================
% Internal APIs Are All You Need  (arXiv:2604.00694)
% Paper I of the trilogy — the wedge. A shared graph of reusable web
% routes: discover once, reuse forever, route the cost through one layer.
% Claims are tagged [shipped] (in the product), [reference] (a runnable,
% tested reference implementation ships in reference/), or [proposed].
% =====================================================================

\usepackage[margin=1in]{geometry}
\usepackage{hyperref}
\usepackage{amsmath}
\usepackage{amssymb}
\usepackage{enumitem}
\usepackage{xcolor}
\hypersetup{colorlinks=true, linkcolor=blue, urlcolor=blue, citecolor=blue}

\newcommand{\impl}{\textcolor{green!55!black}{\small[shipped]}}
\newcommand{\refimpl}{\textcolor{blue!70!black}{\small[reference]}}
\newcommand{\prop}{\textcolor{orange!80!black}{\small[proposed]}}

\title{\textbf{Internal APIs Are All You Need}\\
\large A Shared Route Graph for Agent Web Access:\\
Discover Once, Reuse Forever, Settle over x402}

\author{Lewis Tham\\ Unbrowse AI \\ \texttt{lewis@unbrowse.ai}
  \and Cayden Chik\\ Unbrowse AI \\ \texttt{cayden@unbrowse.ai}}
\date{April 2026 \quad (arXiv:2604.00694)}

\begin{document}
\maketitle

\begin{abstract}
Most AI agents use the web the way a tired human would: open the page, wait for
it to load, fight a popup, click through a menu, wait again. Underneath that
ceremony every site already speaks a structured request language --- the internal
endpoints its own front-end calls. This paper makes one claim and tries not to
oversell it: for the overwhelming majority of agent web tasks, you do not need a
browser per call; you need to learn the internal request path \emph{once} and
reuse it. We frame the web as a shared \emph{route graph}: an intent resolves to a
ranked endpoint shortlist, the chosen route is executed for real data, and the
mapping cost is amortised across every later agent that asks the same question.
The economic core is a single inequality --- reuse wins whenever the fee for
replaying a mapped route undercuts the cost of rediscovering it from a browser ---
and we settle paid execution over \textbf{x402} so the indexer who mapped a route,
the platform, and the site owner are all paid in one on-chain transaction. We
report a peer-reviewed benchmark across 94 live domains and label every claim
\textbf{[shipped]}, \textbf{[reference]}, or \textbf{[proposed]} so the reader is
never sold a roadmap as a result. The short version: internal APIs are a great
first layer --- they are most of what an agent needs to read and act on the web.
What they are \emph{not} is the whole stack, and the companion papers take up
where this one stops.
\end{abstract}

\section{The browser-discovery tax}
The first problem an agent has on the open web is \emph{discovery}: finding the
route in the first place is expensive, and re-deriving it on every call is
wasteful. A browser session loads megabytes of markup, executes the page's
JavaScript, and waits on late XHR settle --- per agent, per call --- to recover a
request the site's own front-end already knew how to make. This is a tax paid over
and over for information that does not change between callers. The thesis of this
paper is that the tax is avoidable: the structured request behind a page is
\emph{learnable once} and \emph{replayable by anyone}, so the discovery cost should
be paid a single time and amortised, not re-paid on every visit.

\section{The shared route graph}
We model the web an agent can act on as a graph of reusable routes. A route is a
callable endpoint with its method, parameters, auth requirements, and a typed
record of what it returned. The two-call contract is deliberately small: an
\emph{intent $\to$ ranked endpoint shortlist $\to$ execute} resolution returns the
candidates that answer an intent, and a single follow-up realizes the chosen
endpoint call against the live site. \impl\ Resolution reads from cache first; a
content-addressed cache stores each intermediate result under the hash of its
content, so the same content resolves identically on any host and a second visit
answers without launching a browser. \impl

The graph is not a private index. A route mapped by one agent is reusable by every
later agent, which is the whole asymmetry the system turns into value: \emph{any
single route is copyable, but the maintained graph is not}. The route graph already
keeps signed records of which routes were used and what they returned. \impl

\section{Why reuse wins: one inequality}
Sharing a route is rational exactly when replaying it is cheaper than rediscovering
it. Let $f_{\text{shared}}$ be the fee an agent pays to reuse a mapped route and
$c_{\text{rediscover}}$ the cost of mapping it from scratch with a browser. The
share-versus-rediscover decision reduces to
\[
  f_{\text{shared}} \;<\; c_{\text{rediscover}},
\]
and because the mapping cost is paid once and amortised over every later caller,
the inequality holds for any route reused more than a handful of times. This is the
adoption condition: a marketplace of routes grows precisely when the fee to reuse
undercuts the cost to re-derive, and the dominant strategy for each agent is to
reuse rather than re-map.

\section{Evaluation: 94 live domains}
We evaluate the API-native path against a browser baseline on a corpus of 94 live
domains spanning search, listings, detail pages, and authenticated reads. On the
shared-route path the system is typically far faster and dramatically cheaper than
driving a browser: a peer-reviewed measurement records a \textbf{3.6$\times$ mean
speedup}, a \textbf{5.4$\times$ median speedup}, and \textbf{40$\times$ fewer
tokens} consumed than the browser baseline across the corpus. The reference
harness reproduces the mechanism in isolation --- a warm cached replay skips the
cold mapping work entirely --- and records a measured $\approx$95$\times$ mean
speedup on the reuse micro-benchmark (\texttt{reference/bench/bench\_reuse.py}),
with a conservative $\approx$27$\times$ observed gap on a single live URL
(\texttt{reference/bench/live-measured-1url.json}). \refimpl\ The release-coverage
methodology --- corpus shape, scoring rubric, and current numbers --- is maintained
alongside the product.

\section{Settlement: one transaction, three parties}
A reused route creates value for three parties, and all three are paid in the same
settled call. Paid execution \emph{settles over \textbf{x402}}: the runtime returns
a payment-required challenge, the caller signs an off-chain authorization, and the
proof rides the retried request. \impl\ The settlement \emph{splits three ways}
--- to the indexer who mapped the route, to the platform, and (when claimed) to the
site owner --- so \emph{x402 routes USDC to all three} in one transaction rather
than a deduction taken after the fact. \impl\ This \emph{three-way live split is
shipped}. \impl\ Discovery and internal-API routing are free; only execution against
a priced route settles, and an \emph{unpaid \texttt{execute} gets HTTP 402}. \impl

\section{Attribution: paid for the value you add}
A maintenance economy must answer a sharper question than ``who touched this
route?'' --- it must answer ``how much better is the route because you touched it?''
We score contribution by the \emph{delta} a maintainer's work makes to a route's
measured quality --- read off freshness, coverage, and faithful-replay between
successive states --- so a no-op edit on top of someone else's work earns nothing
and a real repair is rewarded in proportion to the improvement. This keeps the
incentive pointed at real maintenance rather than at land-grabbing a popular route
and collecting rent on another party's labour. The full accountability machinery
that makes a freshness claim costly to fake --- bonding, challenge, and slashing
over a proof of indexing --- is the subject of the third paper and is not relitigated
here.

\section{What is shipped, what is referenced (no fabricated green)}
In the spirit of not selling a roadmap as a changelog we separate what runs in the
product from what ships as runnable, tested reference code:
\begin{itemize}[leftmargin=1.4em]
  \item \impl\ The two-call contract: \emph{intent $\to$ ranked endpoint shortlist
        $\to$ execute}, resolving cache-first and realizing the chosen endpoint
        call against the live site.
  \item \impl\ The three-tier x402 settlement: every paid execution \emph{splits
        three ways} to indexer, platform, and (when claimed) owner in one
        transaction.
  \item \refimpl\ The reuse benchmark mechanism (\texttt{reference/bench/}): a warm
        replay categorically skips the cold mapping path; the live head-to-head is
        recorded evidence, not a gated claim.
  \item \refimpl\ The signed, content-addressed route ledger
        (\texttt{reference/ledger/}): value off-chain, commitment hash-chained.
\end{itemize}
The line between a running reference and a deployed product is named on each item;
the gap between running code and a press release is the entire difference between a
whitepaper and a pitch, and we stay on the correct side of it.

\section{Conclusion: the first layer, not the last}
The graph is the asset. Every site an agent automates today is re-discovered from
scratch, per agent, per call --- the browser tax paid over and over. Turning that
one-off rediscovery into a shared, reusable route graph --- discover once, reuse
forever, route the cost through one layer instead of re-paying it everywhere ---
removes the tax and makes maintenance, not discovery, the place value compounds.
Internal APIs are most of what an agent needs to read and act on the web, and for
that majority they are, in the affectionate sense of the title, all you need. They
are also not the whole stack: the instant an agent must click, run a shell command,
open a real browser, or put bytes on a socket a bot-detector is watching, the work
leaves the layer where JSON politely comes out. Securing that descent is the second
paper; keeping the shared graph fresh and trustworthy is the third. This one is the
wedge they stand on.

\begin{thebibliography}{9}
\bibitem{security} L.~Tham. \emph{Security Was All You Needed: A Signed,
Layer-Descending Stack for Agent Computation}. Unbrowse AI, 2026.
\bibitem{maintenance} L.~Tham. \emph{Unbrowse Maintenance Network: Proof of
Indexing and Bonded Accountability in a Shared Route Graph}. Unbrowse AI, 2026.
\bibitem{x402} Coinbase. \emph{x402: An Open Protocol for Internet-Native Payments
over HTTP 402}. x402.org whitepaper, 2025. \url{https://github.com/coinbase/x402}.
\end{thebibliography}

\end{document}
